\documentclass[conference]{IEEEtran}
\usepackage{graphicx}
\usepackage{amsmath}
\usepackage{booktabs}
\usepackage{url}

% ============================================================================
% EyeCamXR: The Human Eye as a Camera in Extended Reality
% Skeleton — populate results sections ONLY from measured session artifacts
% (session.json / calibration.json). Novelty guardrails from the 2026-09
% literature review are embedded as comments; do not delete them.
% ============================================================================

\title{EyeCamXR: The Human Eye as a Camera in Shared Extended Reality}

\author{\IEEEauthorblockN{Alexander Vicol and Steve Mann}
\IEEEauthorblockA{Department of Electrical and Computer Engineering,
University of Toronto, Toronto, Canada\\
alexander.vicol@mail.utoronto.ca, mann@eecg.toronto.edu}}

\begin{document}
\maketitle

\begin{abstract}
% TODO after live sessions. Structure:
% (1) We port SSVEP-based reconstruction of viewed imagery (Mann et al.,
%     HealthCom 2019) into a standalone VR headset (Meta Quest 3S) with
%     consumer EEG (Muse 2 / Muse S / Muse S Athena over native BLE);
% (2) frame-exact flicker at 120 Hz with delivered-frequency verification;
% (3) a per-subject calibration gate that measures SSVEP contrast (and the
%     auxiliary-Oz vs.\ behind-the-ear TP9/TP10 gain) before scanning;
% (4) shared-XR spectation: the image reconstructs progressively on
%     spectators' screens/headsets as the participant's brain responds;
% (5) results: phantom-validated pipeline (r = X), live sessions (r = Y).
\end{abstract}

\section{Introduction}
% Lineage: Mann et al., "The Human Eye as a Camera," IEEE HealthCom 2019 --
% flickering stimulus at 12/15 Hz, SSVEP relative power (15+30 Hz over
% 14-50 Hz) at each gaze position = pixel; auxiliary Oz electrode on a Muse.
% NOVELTY GUARDRAILS (verified 2026-09-11):
%  - SAFE: first SSVEP imaging inside an HMD; first Quest-family SSVEP study.
%  - NOT novel: harmonic fusion (Garcia, Zheng & Mann, IEEE NER 2021);
%    lock-in detection (Garcia et al., SIBGRAPI 2020) -- cite both.
%  - MUST cite and distinguish: Wang, Huang, Muckli, Faccio, npj Biomedical
%    Innovations 2026 (imagined-image drawing via frequency-multiplexed
%    SSVEP, single Oz channel; imagined, not viewed; does not cite Mann).
%  - Distinguish from learned generative EEG-to-image decoding (e.g.
%    arXiv:2308.13234, Alljoined1): ours is a per-pixel physical measurement,
%    not a learned prior.

\section{Background and Related Work}
% - SSVEP in HMDs: Koo et al. EMBC 2015 (VR-HMD SSVEP viable, ITR +10%);
%   Frontiers Neurosci 2024 (PICO Neo3, black-frame insertion degrades
%   spectral purity); Frontiers Neurosci 2026 (Vision Pro, f = 90/n
%   constraint; MR vs VR SNR 1.43 vs 0.58 dB; accuracy falls with depth).
% - Frequency mislabeling costs 20+ pp (XR timing compensation, HoloLens/
%   Moverio study) -> motivates our delivered-frequency measurement.
% - Behind-the-ear vs occipital SSVEP: Chang et al., Sensors 18(2):615, 2018.
% - Muse aux input: 4 EEG + 2 amplified aux channels (Muse S spec);
%   Athena 8-channel presets (empirical channel count reported here).

\section{System}
\subsection{Architecture}
% Quest 3S Unity app (all-in-one): native libmuse BLE (model-aware presets:
% PRESET_20/21 classic, PRESET_1021/1022 Athena; runtime EegChannelCount
% probe) OR OSC from the MuseLog phone app; on-device DSP is a bin-exact
% port of the reference pipeline (golden-vector parity <= 1e-9 vs
% scipy.signal.welch); spectator bridge (UDP JSON -> WebSocket) streams the
% progressive reconstruction to any browser, including a second headset.
\subsection{Stimulus}
% Region-serial guided cursor (whole-image flicker fails: peripheral SSVEP
% contamination, Mann et al. 2019 Sec. III). 120 Hz refresh: 4 frames on /
% 4 off = exactly 15 Hz; 12 Hz (10-frame period) and the 30 Hz harmonic are
% also frame-exact -- the only Quest 3S rate with this property.
% Delivered flicker frequency is re-measured from logged frame timestamps;
% sessions gate on |measured - nominal| <= 0.1 Hz and dropped frames < 1%.
\subsection{Calibration and artifact handling}
% 6 x (7 s ON / 7 s OFF) blocks; per-channel robust d' on relative SSVEP
% power; scan refused unless best d' >= 1.0, on/off ratio >= 1.3, rank
% >= 5/6. Weights = normalized (ratio - 1) over admitted channels.
% Per-segment median baseline + MAD-floored sigma clipping; cell-channels
% with > 20% clipped samples dropped (blink/motion).

\section{Validation Without a Subject}
% Phantom subject: real-time synthetic observer that tails the live stimulus
% logs and injects SSVEP at the MEASURED flicker rate over real UDP, with
% asymmetric channel SNR (temporal >> frontal) and minutes-scale electrode
% drift. Laptop reference: full-stack gate r = 0.849. Unity port: Play-Mode
% phantom gate r = [TODO]; on-device phantom-at-Quest r = [TODO].

\section{Live Results}
% TODO: n subjects, TP9/TP10-only vs aux-Oz d' comparison table (this
% comparison on a Muse is itself unreported in the literature), example
% reconstructions with r, scan durations per preset.

\section{Shared XR Spectation}
% Progressive row-by-row build-up on spectators' displays; late-join replay;
% message schema versioned for a native spectator headset app.

\section{Limitations and Future Work}
% - LCD low-persistence strobing: flicker is an amplitude envelope on the
%   display's own strobe; we report measured spectra, photodiode validation
%   is future work.
% - Color: sequential per-color-calibrated R/G/B passes (chromatic SSVEP
%   gain differs per color and interacts with frequency -- Sensors 2018);
%   frequency-multiplexed RGB (12/15/20 Hz, harmonics non-colliding at
%   120 Hz) as future single-pass mode.
% - Audio ("ear as a microphone"): 40 Hz AM-tone ASSR is minutes-scale at
%   useful SNR; a Muse ASSR validation would itself be novel. Framed as a
%   slow envelope demonstration, never real-time audio.
% - Chirplet-transform VEP estimation (ACT) as the next detector.

\section{Ethics and Data Sharing}
% REB status; recordings exportable as CSV (MuseLog-compatible schema);
% no audio/video of bystanders; participant consent for shared spectation.

\bibliographystyle{IEEEtran}
\bibliography{references}

\end{document}
